\uuid{jgvD}
\titre{Distance}
\auteur{bodin ; hanani ; mzari}
\organisation{exo7}
\licence{CC BY-NC-SA 4.0}
\datecreate{2020-01-20}
\niveau{L1}
\module{Géométrie}
\chapitre{Géométrie du plan}
\sousChapitre{Distance, norme, produit scalaire}
\difficulte{2}
\theme{distance}
\exo7qcmid{lille-1.geometrie-plane.018}

\contenu{
\begin{question}
\qclassification{140.01,140.02}
\qtags{motcle=distance}
Dans le plan rapporté à un repère orthonormé, on considère les points $\displaystyle A(1,2)$, $B(2,1)$ et $C(-2,1)$. Quelle est la distance $d$ du point $C$ à la droite $(AB)$ ?
\begin{answers}
    \bad{$d=\sqrt{2}$}
    \bad{$d=3$}
    \good{$d=2\sqrt{2}$}
    \bad{$d=\sqrt{10}$}
\end{answers}
\begin{explanations}
Utiliser la formule $\displaystyle d=\frac{\left|\mbox{det}\left(\overrightarrow{AC},\overrightarrow{AB}\right)\right|}{\|\overrightarrow{AB}\|}=2\sqrt{2}$.
\end{explanations}
\end{question}
}
